\documentclass{article}
\usepackage{annot}
\begin{document}

% ============================================================================
% Chapitre 4 — Comparaison de deux populations
% (pages blanches insérées après les diapos 11, 15, 27, 34 et 40)
% ============================================================================

% --- Où en sommes-nous ? ----------------------------------------------------
\annot{3}{Où en sommes-nous}{
  \entourer{(64.3,73.2)}{3.2}{2.4}
  \entourer{(79.5,61.5)}{13.5}{2.8}
}

% --- Les 5 situations -------------------------------------------------------
\annot{5}{Les 5 situations}{
  \entourer[rouge]{(86,24.7)}{5.5}{2.1}
  \ecrire[violet][7]{(129,57.5)}{MÊME UNITÉ\\MESURÉE 2 FOIS}
  \ecrire[vert][10]{(12,13)}{EX. ZINC : TEST $F$ (QUESTION 1), PUIS TEST $t$ (QUESTION 2)}
}

% --- Feuille de route -------------------------------------------------------
\annot{6}{Feuille de route}{
  \ecrire[orange][9]{(126,40.5)}{INDÉPENDANTS :\\LE TEST $F$\\DÉCIDE !}
  \fleche[orange][-20]{(127,32)}{(123.5,30.5)}
  \ecrire[violet][9]{(3,20)}{APPARIÉ : MÊME\\SUJET MESURÉ 2 FOIS}
  \fleche[violet][20]{(22,21)}{(33,25)}
}

% --- Exemple 1 : zinc -------------------------------------------------------
\annot{7}{Exemple 1 : Le zinc}{
  \entourer[rouge]{(111.3,61.4)}{8.3}{2.8}
  \ecrire[rouge][9]{(38,59.3)}{$\Rightarrow H_1 : \mu_1 > \mu_2$ (UNILATÉRAL À DROITE)}
  \fleche[rouge][-25]{(104,58.8)}{(102,57.2)}
  \surligner[jaune]{(122,20.3)}{(147,30.2)}
  \ecrire[violet][9]{(119,16.5)}{$\rightarrow$ TEST $F$}
}

\annot{8}{Exemple 1 : Toujours regarder}{
  \ecrire[vert][9]{(8,29)}{MÉDIANES : 5,00 ET 4,50\\
    ÉTENDUES : $6{,}51-3{,}90 = 2{,}61$\\
    \phantom{ÉTENDUES : }$6{,}03-2{,}58 = 3{,}45$\\
    $\bar x_1 - \bar x_2 = 5{,}00 - 4{,}39 = 0{,}61$}
}

\annot{9}{QUIZ}{
  \entourer[rouge]{(44.5,49)}{6.2}{2.4}
  \entourer[rouge]{(92.5,24.7)}{16.5}{2.7}
  \ecrire[violet][9]{(10,13)}{CE QUE LES BIOLOGISTES SOUPÇONNENT VA DANS $H_1$ !\\
    INDÉPENDANTS $\Rightarrow$ TEST $F$ (Q1) PUIS TEST $t$ À 2 ÉCHANTILLONS (Q2)}
}

% --- Loi de Fisher ----------------------------------------------------------
\annot{11}{Situation et distribution}{
  \surligner[rose]{(33,49.3)}{(46,56)}
  \surligner[rose]{(62.5,49.3)}{(75,56)}
  \ecrire[violet][9]{(104,56.5)}{$= W_1$ ET $W_2$,\\INDÉPENDANTES}
  \ecrire[rouge][9]{(107,23.5)}{SOUS $H_0$ ($\sigma_1^2 = \sigma_2^2$) :\\$F = S_1^2 / S_2^2$}
}
\apres{11}{
  \ecrire[violet][13]{(8,84)}{D'OÙ VIENT LA LOI $F$ ?}
  \ecrire[vert][12]{(8,76)}{%
    $W_1 = \dfrac{(n_1-1)S_1^2}{\sigma_1^2} \sim \chi^2_{n_1-1}$ \quad ET \quad
    $W_2 = \dfrac{(n_2-1)S_2^2}{\sigma_2^2} \sim \chi^2_{n_2-1}$ \ (INDÉP.)}
  \ecrire[vert][12]{(8,56)}{%
    $F = \dfrac{W_1/(n_1-1)}{W_2/(n_2-1)} = \dfrac{S_1^2/\sigma_1^2}{S_2^2/\sigma_2^2} \sim \mathcal{F}(n_1-1,\ n_2-1)$}
  \ecrire[rouge][12]{(8,36)}{SOUS $H_0$ : $\sigma_1^2 = \sigma_2^2$ $\Rightarrow$ LES $\sigma^2$ S'ANNULENT :
    $F_{obs} = \dfrac{s_1^2}{s_2^2}$}
  \ecrire[violet][11]{(8,18)}{$F_{obs}$ PRÈS DE 1 : VARIANCES SEMBLABLES\\
    $F_{obs}$ TRÈS PETIT OU TRÈS GRAND : ON REJETTE $H_0$}
  % petite loi F asymétrique avec deux queues
  \begin{scope}[shift={(112,40)}, x=5.5mm, y=20mm]
    \fill[orange, opacity=.45] (0.02,0) -- plot[domain=0.02:0.228, samples=10, smooth] (\x, {34.02*\x*\x*(1+1.2*\x)^(-5.5)}) -- (0.228,0) -- cycle;
    \fill[orange, opacity=.45] (4.95,0) -- plot[domain=4.95:7.5, samples=15, smooth] (\x, {34.02*\x*\x*(1+1.2*\x)^(-5.5)}) -- (7.5,0) -- cycle;
    \draw[crayon lisse, bleu] plot[domain=0.01:7.5, samples=60, smooth] (\x, {34.02*\x*\x*(1+1.2*\x)^(-5.5)});
    \draw[crayon lisse, bleu] (0,0) -- (7.6,0);
  \end{scope}
  \ecrire[bleu][10]{(110,61.5)}{ASYMÉTRIQUE, $F > 0$}
}

\annot{12}{Propriétés de la loi de Fisher}{
  \entourer[rouge]{(63.5,42)}{15}{2.6}
  \ecrire[violet][9]{(64,30.5)}{$\nu_1 \leftrightarrow \nu_2$\\INVERSÉS !}
  \ecrire[vert][9]{(61,10)}{$F_{0,95;\,6,\,5} = \frac{1}{F_{0,05;\,5,\,6}} = \frac{1}{4{,}39} = 0{,}228$}
}

\annot{13}{Intervalle de confiance pour}{
  \entourer[rouge]{(100,43.5)}{9.5}{3.2}
  \ecrire[vert][9]{(12,38)}{EX. ZINC (90 \%) : $\left[\frac{0{,}606}{4{,}95}\ ;\ 0{,}606 \times 4{,}39\right] = [0{,}12\ ;\ 2{,}66]$ \quad CONTIENT 1}
  \ecrire[rouge][9]{(106,18.5)}{1 $\in$ IC $\Leftrightarrow$ ON NE\\REJETTE PAS $H_0$\\(BILATÉRAL, SEUIL $\alpha$)}
}

\annot{14}{Test de Fisher pour}{
  \ecrire[violet][9]{(52,64.5)}{$\leftarrow$ DL DU NUMÉRATEUR, DL DU DÉNOMINATEUR}
  \ecrire[violet][8]{(103,35.8)}{BILATÉRAL : $\alpha/2$ DE CHAQUE CÔTÉ}
}

% --- Exemple 1 : test de Fisher -------------------------------------------
\annot{15}{Exemple 1 (suite) : Test de Fisher}{
  \entourer[violet]{(116.3,42.2)}{4}{2.5}
  \ecrire[violet][9]{(120,58)}{DL INVERSÉS (5, 6)\\POUR LA BORNE\\INFÉRIEURE}
  \fleche[violet][-20]{(128,46)}{(119.5,44.5)}
}
\apres{15}{
  \ecrire[violet][13]{(8,84)}{BORNE INFÉRIEURE : $c = F_{0,95;\,6,\,5}$ = ?}
  \ecrire[vert][12]{(8,72)}{%
    $P(F_{6,5} < c) = 0{,}05$\\[3mm]
    $\iff P\!\left(\dfrac{1}{F_{6,5}} > \dfrac{1}{c}\right) = 0{,}05$\\[3mm]
    OR $\dfrac{1}{F_{6,5}} \sim F_{5,6}$ \ DONC \ $\dfrac{1}{c} = F_{0,05;\,5,\,6} = 4{,}39$ (TABLE)\\[3mm]
    $\Rightarrow c = \dfrac{1}{4{,}39} = 0{,}228$}
  \encadrer[rouge]{(5,23.5)}{(52,37.5)}
  % loi F(6,5) : deux queues de 5 %
  \begin{scope}[shift={(104,12)}, x=8.5mm, y=32mm]
    \fill[orange, opacity=.45] (0.02,0) -- plot[domain=0.02:0.228, samples=10, smooth] (\x, {34.02*\x*\x*(1+1.2*\x)^(-5.5)}) -- (0.228,0) -- cycle;
    \fill[orange, opacity=.45] (4.95,0) -- plot[domain=4.95:6.2, samples=15, smooth] (\x, {34.02*\x*\x*(1+1.2*\x)^(-5.5)}) -- (6.2,0) -- cycle;
    \draw[crayon lisse, bleu] plot[domain=0.01:6.2, samples=60, smooth] (\x, {34.02*\x*\x*(1+1.2*\x)^(-5.5)});
    \draw[crayon lisse, bleu] (0,0) -- (6.3,0);
    \draw[crayon lisse, vert, line width=1.1pt] (0.606,0) -- (0.606,0.72);
  \end{scope}
  \ecrire[orange][10]{(88,20)}{0,05}
  \fleche[orange][20]{(94,15)}{(104.8,12.8)}
  \ecrire[orange][10]{(138,26)}{0,05}
  \fleche[orange][-20]{(144,20)}{(149,13.3)}
  \ecrirec[rouge][10]{(106,7.5)}{0,228}
  \ecrirec[rouge][10]{(146,7.5)}{4,95}
  \ecrire[vert][10]{(118,40)}{$F_{obs} = 0{,}606$}
  \fleche[vert][15]{(119,35.5)}{(110,34)}
  \ecrirec[violet][10]{(126,7.5)}{NON-REJET}
}

\annot{16}{Exemple 1 (suite) : La décision}{
  \draw[crayon lisse, vert, pointe] (75,23.5) -- (43,23.5);
  \draw[crayon lisse, vert, pointe] (75,23.5) -- (106.3,23.5);
  \ecrirec[vert][8]{(65,26.4)}{NON-REJET (AIRE 0,90)}
  \ecrire[vert][8]{(112,67.5)}{VALEUR-P $= 2 \times P(F_{6,5} < 0{,}606)$\\$= 0{,}557$}
}

\annot{17}{Exemple 1 (suite) : Sortie R de}{
  \surligner[jaune]{(21,47.8)}{(40,51.2)}
  \surligner[rose]{(40.5,47.8)}{(69.5,51.2)}
  \surligner[jaune]{(70.5,47.8)}{(98.5,51.2)}
  \surligner[jaune]{(21.5,39.3)}{(50,42.6)}
  \ecrire[rouge][9]{(100,54.5)}{$0{,}557 > 0{,}10$ : ON NE REJETTE PAS $H_0$}
  \ecrire[vert][9]{(116,45)}{IC À 90 \% :\\$[0{,}12\ ;\ 2{,}66]$ CONTIENT 1}
  \fleche[vert][10]{(115,40)}{(51.5,41)}
}

\annot{18}{QUIZ}{
  \entourer[rouge]{(94.5,53.6)}{14}{2.6}
  \ecrire[violet][10]{(18,11)}{$1{,}652 \approx 1/0{,}606$ ET $4{,}39 = 1/0{,}228$ : C'EST LE MÊME TEST, VU À L'ENVERS}
}

\annot{19}{Attention}{
  \entourer[rouge]{(110,57.7)}{9.5}{2.6}
  \ecrire[rouge][11]{(122,63)}{!!}
}

% --- Deux moyennes, variances égales ---------------------------------------
\annot{21}{Cas 1 : Variances égales}{
  \ecrire[vert][10]{(109,51)}{SI $n_1 = n_2$ :\\[0.5mm]$S_p^2 = \dfrac{S_1^2 + S_2^2}{2}$}
  \ecrire[violet][9]{(72,9)}{$\leftarrow$ PLUS PETITE QUAND $n_1$ ET $n_2$ SONT GRANDS}
}

\annot{22}{Intervalle de confiance pour}{
  \ecrire[violet][9]{(13,37.5)}{ESTIMATION\\PONCTUELLE $\rightarrow$}
  \ecrire[rouge][9]{(114,38.5)}{$\leftarrow$ MARGE\\D'ERREUR}
  \ecrire[vert][9]{(118,64.5)}{ZINC :\\$\nu = 7+6-2 = 11$}
}

\annot{23}{Test d'hypothèses pour}{
  \ecrire[violet][9]{(101,55.5)}{$= \dfrac{\text{ESTIMATION} - 0}{\widehat{\text{ET}}}$\\[1mm]($\mu_1 - \mu_2 = 0$ SOUS $H_0$)}
}

\annot{24}{Exemple 1 (suite) : Comparaison}{
  \ecrire[vert][9]{(122,50)}{VALEUR-P :\\$P(T_{11} > 1{,}01)$\\TABLE :\\$t_{11;\,0,15} = 1{,}088$\\$t_{11;\,0,20} = 0{,}876$\\$\Rightarrow\ ]0{,}15\ ;\ 0{,}20[$}
}

\annot{25}{Exemple 1 (suite) : Région critique}{
  \ecrire[rouge][9]{(122,68.5)}{$t_{obs}$ HORS DE LA\\RÉGION CRITIQUE\\$\Leftrightarrow$ VALEUR-P $> \alpha$}
}

\annot{26}{Exemple 1 (suite) : Sortie R de}{
  \entourer[vert]{(62.5,74.4)}{18}{2.5}
  \surligner[jaune]{(16,56.8)}{(31.5,60.2)}
  \surligner[rose]{(32,56.8)}{(41.6,60.2)}
  \surligner[jaune]{(42.2,56.8)}{(67,60.2)}
  \surligner[jaune]{(16,47.9)}{(44.2,51.3)}
  \ecrire[vert][9]{(72,51.8)}{$H_1 : \mu_1 > \mu_2$ $\Rightarrow$ IC UNILATÉRAL $[-0{,}475\ ;\ \infty[$\\$0 \in$ IC $\Rightarrow$ ON NE REJETTE PAS $H_0$}
}

\annot{27}{Exemple 1 (suite) : Test bilatéral}{
  \entourer[vert]{(62,62.7)}{16.5}{2.3}
  \surligner[jaune]{(42.2,47.8)}{(67,51.2)}
  \surligner[jaune]{(16,39.7)}{(46.5,43)}
  \ecrire[rouge][7.5]{(69,51.6)}{$= 2 \times 0{,}1672$}
  \ecrire[vert][9]{(49,43.6)}{$\leftarrow$ CONTIENT 0 (CALCUL À LA MAIN : PAGE SUIVANTE)}
}
\apres{27}{
  \ecrire[violet][13]{(8,84)}{IC À 95 \% POUR $\mu_1 - \mu_2$ (ZINC)}
  \ecrire[vert][12]{(8,73)}{%
    $(\bar x_1 - \bar x_2) \pm t_{11;\,0,025}\ s_p\sqrt{\frac17 + \frac16}$\\[3mm]
    $= 0{,}61 \pm 2{,}201 \times 0{,}604$\\[3mm]
    $= 0{,}61 \pm 1{,}33$\\[3mm]
    $= [-0{,}72\ ;\ 1{,}94]$ \ $\mu$G/L}
  \encadrer[rouge]{(6,38)}{(64,46.5)}
  \ecrire[violet][11]{(80,72)}{MARGE D'ERREUR :\\$2{,}201 \times 0{,}604 = 1{,}33$}
  % droite des valeurs
  \draw[crayon lisse, vert] (10,16) -- (150,16);
  \draw[crayon lisse, rouge, line width=1.4pt] (40,16) -- (126,16);
  \draw[crayon lisse, rouge, line width=1.2pt] (42,13) -- (40,13) -- (40,19) -- (42,19);
  \draw[crayon lisse, rouge, line width=1.2pt] (124,13) -- (126,13) -- (126,19) -- (124,19);
  \draw[crayon lisse, violet, line width=1.2pt] (64,12.5) -- (64,19.5);
  \fill[vert] (85,16) circle (1mm);
  \ecrirec[rouge][10]{(40,9)}{$-0{,}72$}
  \ecrirec[rouge][10]{(126,9)}{$1{,}94$}
  \ecrirec[violet][11]{(64,23)}{0}
  \ecrirec[vert][10]{(85,9)}{$0{,}61$}
  \ecrire[rouge][11]{(80,52)}{$0 \in$ IC $\Rightarrow$ ON NE REJETTE PAS\\$H_0 : \mu_1 = \mu_2$ (BILATÉRAL, 5 \%)\\[1mm]
    \textcolor{vert}{R : $[-0{,}720\ ;\ 1{,}940]$ \ OK !}}
}

\annot{28}{QUIZ}{
  \draw[crayon lisse, vert] (28,38.5) -- (142,38.5);
  \draw[crayon lisse, rouge, line width=1.2pt] (41.2,37) -- (39.6,37) -- (39.6,40) -- (41.2,40);
  \draw[crayon lisse, rouge, line width=1.2pt] (129.5,37) -- (131.1,37) -- (131.1,40) -- (129.5,40);
  \draw[crayon lisse, violet, line width=1.1pt] (64.4,36.8) -- (64.4,40.2);
  \fill[vert] (85.8,38.5) circle (0.8mm);
  \ecrirec[rouge][8]{(39.6,42.2)}{$-0{,}72$}
  \ecrirec[violet][8]{(64.4,42.2)}{0}
  \ecrirec[vert][8]{(85.8,42.2)}{0,61}
  \ecrirec[rouge][8]{(131.1,42.2)}{1,94}
  \ecrire[rouge][10]{(18,10)}{0 EST PLAUSIBLE... MAIS UNE DIFFÉRENCE DE 1,9 AUSSI !}
}

% --- Deux moyennes, variances inégales -------------------------------------
\annot{29}{Cas 2 : Variances inégales}{
  \ecrire[rouge][9]{(113,51.5)}{$\leftarrow$ PAS DE $S_p$ :\\CHAQUE VARIANCE\\SÉPARÉMENT}
  \ecrire[violet][9]{(108,28.5)}{$\min(n_1-1,\,n_2-1) \le \nu$\\$\le n_1+n_2-2$\\EX. 2 : $7{,}48 \rightarrow 7$}
}

\annot{30}{Cas 2 : Règles de décision}{
  \ecrire[vert][8]{(116,76)}{ZINC ($t_{obs}$ ; $\nu$ ; VALEUR-P) :\\
    $S_p$ : \ 1,01 ; 11 ; 0,167\\
    WELCH : 0,99 ; 9,42 ; 0,174}
  \ecrire[rouge][9]{(115,38.5)}{$\nu$ : WELCH\\(ARRONDI VERS\\LE BAS)}
}

\annot{31}{Exemple 2 : Mercure}{
  \entourer[rouge]{(142,50.4)}{5.8}{4.6}
  \ecrire[rouge][8]{(100,46.2)}{$0{,}2524 / 0{,}0521 \approx 4{,}8$ !}
}

\annot{32}{Exemple 2 (suite) : Étape 1}{
  \ecrire[violet][8.5]{(74,47.8)}{$\alpha/2 = 0{,}05$ ; $\nu_1 = 8-1 = 7$ (NUM.), $\nu_2 = 10-1 = 9$ (DÉN.)}
  \ecrire[rouge][9]{(116,41)}{$23{,}46 \gg 3{,}29$ :\\TRÈS LOIN !}
  \ecrire[vert][9]{(18,13)}{AVEC $s_2^2$ EN HAUT : $F_{obs} = 1/23{,}46 = 0{,}043 < 1/F_{0,05;\,9,\,7} = 1/3{,}68 = 0{,}27$\\
    $\Rightarrow$ MÊME CONCLUSION}
}

\annot{33}{Exemple 2 (suite) : Étape 2}{
  \ecrire[rouge][9]{(118,47)}{$\leftarrow$ ARRONDI\\VERS LE BAS\\(PRUDENT)}
  \ecrire[vert][8.5]{(95,34)}{TABLE : $t_{7;\,0,005} = 3{,}499 < 3{,}90$\\$\Rightarrow$ VALEUR-P (BILAT.) $< 0{,}01$}
}

\annot{34}{Exemple 2 (suite) : Sortie R}{
  \surligner[jaune]{(9.9,48.7)}{(26,52)}
  \surligner[rose]{(27,48.7)}{(46,52)}
  \surligner[jaune]{(46.5,48.7)}{(88,52)}
  \surligner[jaune]{(10.5,38.8)}{(46.5,42.2)}
  \ecrire[vert][9]{(50,43.3)}{$\leftarrow$ À LA MAIN ($\nu = 7$) : $[0{,}139\ ;\ 0{,}569]$ (PAGE SUIVANTE)}
}
\apres{34}{
  \ecrire[violet][13]{(8,84)}{IC À 95 \% POUR $\mu_1 - \mu_2$ (MERCURE, WELCH)}
  \ecrire[vert][12]{(8,73)}{%
    $(\bar x_1 - \bar x_2) \pm t_{7;\,0,025}\ \sqrt{\frac{s_1^2}{n_1} + \frac{s_2^2}{n_2}}$\\[3mm]
    $= 0{,}354 \pm 2{,}365 \times 0{,}0908$\\[3mm]
    $= 0{,}354 \pm 0{,}215$\\[3mm]
    $= [0{,}139\ ;\ 0{,}569]$ \ $\mu$G/G}
  \encadrer[rouge]{(6,38)}{(66,46.5)}
  \ecrire[rouge][11]{(82,72)}{$0 \notin$ IC $\Rightarrow$ ON REJETTE $H_0$\\(COHÉRENT AVEC LE TEST)}
  \ecrire[violet][11]{(82,52)}{R ($\nu = 7{,}48$) : $[0{,}142\ ;\ 0{,}566]$\\
    UN PEU PLUS ÉTROIT CAR\\$t_{7,48;\,0,025} = 2{,}33 < 2{,}365$}
  \ecrire[orange][11]{(8,22)}{ARRONDIR $\nu$ VERS LE BAS $\Rightarrow$ $t$ PLUS GRAND $\Rightarrow$ IC PLUS LARGE (PRUDENT)}
}

\annot{35}{QUIZ}{
  \ecrire[violet][9.5]{(18,12)}{AVEC $S_p$ : $\widehat{\text{ET}} = 0{,}171\sqrt{\frac18 + \frac1{10}} = 0{,}081$ AU LIEU DE $0{,}091$ (WELCH) $\Rightarrow$ $t_{obs}$ GONFLÉ}
}

% --- Données appariées -----------------------------------------------------
\annot{37}{Deux scénarios}{
  \ecrire[rouge][10]{(30,6)}{$\rightarrow$ SCÉNARIO 2 (APPARIÉ) !}
}

\annot{38}{Notion d'appariement}{
  \surligner[jaune]{(80,60.1)}{(116.5,64.1)}
  \surligner[jaune]{(79,47)}{(103.5,51)}
  \ecrire[violet][10]{(18,13.5)}{APPARIER = COMPARER A ET B À L'INTÉRIEUR D'UNE MÊME UNITÉ\\(FERME, PATIENT, SITE, FEUILLE...)}
}

\annot{39}{Notion de données appariées}{
  \surligner[jaune]{(102,60.8)}{(146.5,64.8)}
  \surligner[jaune]{(14,55.7)}{(34,59.7)}
}

\annot{40}{Méthode d'analyse}{
  \ecrire[rouge][9]{(119,40.5)}{CHAP. 3C :\\[0.5mm]$t = \dfrac{\bar X - \mu_0}{S/\sqrt n}$\\[0.5mm]AVEC $\mu_0 = 0$}
}
\apres{40}{
  \ecrire[violet][13]{(8,84)}{POURQUOI L'APPARIEMENT AIDE-T-IL ?}
  \ecrire[vert][12]{(8,73)}{%
    $Var(D) = Var(X_1 - X_2) = \sigma_1^2 + \sigma_2^2 - 2\,Cov(X_1, X_2)$}
  \ecrire[rouge][12]{(8,60)}{MÊME SUJET $\Rightarrow$ $Cov(X_1,X_2) > 0$ $\Rightarrow$ $Var(D)$ BEAUCOUP PLUS PETITE}
  \ecrire[vert][11]{(8,47)}{EX. 3 :\quad $s_{\text{AVANT}}^2 + s_{\text{APRÈS}}^2 = 11{,}6^2 + 9{,}8^2 \approx 230$\\[2mm]
    \phantom{EX. 3 :}\quad MAIS \ $s_D^2 = 3{,}16^2 \approx 10$ \ ! \quad (CORRÉLATION $\approx 0{,}97$)}
  % petit tableau des paires
  \draw[crayon, violet] (20,18.5) -- (102,18.5);
  \draw[crayon, violet] (40,4) -- (40,24);
  \draw[crayon, violet] (60,4) -- (60,24);
  \draw[crayon, violet] (80,4) -- (80,24);
  \ecrirec[violet][10]{(30,22)}{PAIRE}
  \ecrirec[violet][10]{(50,22)}{$X_1$}
  \ecrirec[violet][10]{(70,22)}{$X_2$}
  \ecrirec[rouge][10]{(91,22)}{$D = X_1 - X_2$}
  \ecrirec[violet][11]{(30,14.5)}{1}
  \ecrirec[violet][10]{(50,14.5)}{$X_{11}$}
  \ecrirec[violet][10]{(70,14.5)}{$X_{21}$}
  \ecrirec[rouge][10]{(91,14.5)}{$D_1$}
  \ecrirec[violet][10]{(30,8)}{$\vdots$}
  \ecrirec[rouge][10]{(91,8)}{$\vdots$}
  \ecrire[orange][11]{(108,20)}{$\rightarrow$ UN SEUL\\ÉCHANTILLON\\DE $n$ DIFFÉRENCES}
}

\annot{41}{Test apparié : règles}{
  \ecrire[violet][9.5]{(30,47.5)}{= TEST $t$ À UN ÉCHANTILLON SUR LES $d_i$ ($\mu_0 = 0$), \ DL : $n - 1$}
}

\annot{42}{Exemple 3 : Un antihypertenseur}{
  \surligner[jaune]{(24,41.3)}{(148,45.5)}
  \entourer[rouge]{(103.8,16.6)}{4.5}{2.7}
  \ecrire[rouge][10]{(18,10)}{$s_D \ll 11{,}6$ ET $9{,}8$ : ON A RETIRÉ LA VARIABILITÉ ENTRE PATIENTS}
}

\annot{43}{Exemple 3 (suite) : L'appariement}{
  \ecrire[violet][9]{(66,56)}{CHAQUE PATIENT\\EST SON PROPRE\\TÉMOIN $\rightarrow$}
  \ecrire[rouge][9]{(66,38)}{ON ANALYSE\\LES PENTES\\(LES $d_i$)}
}

\annot{44}{Exemple 3 (suite) : Solution}{
  \ecrire[vert][8.5]{(116,70.5)}{IC 95 \% POUR $\mu_D$ :\\$8{,}3 \pm 2{,}262 \times 1{,}001$\\$= 8{,}3 \pm 2{,}26$\\$= [6{,}04\ ;\ 10{,}56]$\\BAISSE : 6 À 10,6 MMHG}
}

\annot{45}{QUIZ}{
  \entourer[rouge]{(28.5,32.8)}{10.5}{2.8}
  \ecrire[violet][10]{(18,17.5)}{APPARIÉ : $\widehat{\text{ET}} = 3{,}164/\sqrt{10} = 1{,}00$\\
    INDÉPENDANTS (À TORT) : $\widehat{\text{ET}} = 10{,}72\sqrt{2/10} = 4{,}80$ !}
}

\annot{46}{Lecture d'une sortie R}{
  \entourer[vert]{(82,63.7)}{15}{2.3}
  \surligner[jaune]{(17.8,48.1)}{(34.3,51.5)}
  \surligner[rose]{(35.5,48.1)}{(44,51.5)}
  \surligner[jaune]{(45,48.1)}{(70.5,51.5)}
  \surligner[jaune]{(37,31.4)}{(50,34.8)}
  \ecrire[vert][9]{(112,58)}{$H_1 : \mu_D > 0$\\(FOND > SURFACE)\\[1mm]$t_{5;\,0,01} = 3{,}365$\\$3{,}69 > 3{,}365$\\$\Rightarrow$ ON REJETTE $H_0$}
}

\annot{47}{Quelques remarques}{
  \surligner[jaune]{(31.5,70.3)}{(41.5,74.3)}
}

\annot{48}{QUIZ}{
  \souligner[violet]{(18,63)}{(67,63)}
  \souligner[violet]{(122.5,49.4)}{(131.5,49.4)}
  \souligner[violet]{(76,30.8)}{(126,30.8)}
}

% --- Deux proportions -------------------------------------------------------
\annot{50}{Comparaison de deux proportions}{
  \ecrire[rouge][9]{(47,38.5)}{$\leftarrow = \dfrac{\text{NB TOTAL DE SUCCÈS}}{\text{NB TOTAL D'ESSAIS}}$ \ (SOUS $H_0$ : UN SEUL $p$)}
  \ecrire[violet][9]{(108,21)}{$= \dfrac{\text{ESTIMATION} - 0}{\widehat{\text{ET}} \text{ SOUS } H_0}$}
}

\annot{51}{Règles de décision, conditions}{
  \surligner[jaune]{(18,62)}{(33,66)}
  \ecrire[violet][9]{(18,10)}{SOUS $H_0$, $p_1 = p_2$ : ON COMBINE TOUTE L'INFORMATION POUR ESTIMER\\CETTE PROPORTION COMMUNE}
}

\annot{52}{Exemple 4 : Taux de guérison}{
  \ecrire[vert][9]{(132,56.5)}{$= 90/120$\\$= 66/110$}
  \ecrire[rouge][10]{(92,43)}{$\Rightarrow$ BILATÉRAL ($\ne$)}
  \coche{(64,4.5)}
}

\annot{53}{Exemple 4 (suite) : Solution}{
  \surligner[rose]{(22.5,64.3)}{(42,68.5)}
  \surligner[jaune]{(80.5,32)}{(101,39.5)}
  \ecrire[violet][9]{(62,74)}{$\leftarrow$ $\widehat{P}$ COMBINÉE}
  \ecrire[vert][9]{(119,72.5)}{VALEUR-P :\\$P(Z > 2{,}43)$\\$= 1 - 0{,}9925$\\$= 0{,}0075$\\$\times 2 = 0{,}015$}
}

\annot{54}{Exemple 4 (suite) : La décision}{
  \draw[crayon lisse, vert, line width=1pt] (49.1,36.3) -- (49.1,41);
  \ecrirec[vert][8]{(50.5,45)}{$-2{,}43$}
  \ecrire[vert][8.5]{(5,71.5)}{VALEUR-P\\$= 2 \times P(Z > 2{,}43)$\\$= 2 \times 0{,}0075$\\$= 0{,}015 < 0{,}05$}
}

\annot{55}{Exemple 4 (suite) : Sortie R}{
  \surligner[jaune]{(10,56.2)}{(34.5,59.6)}
  \surligner[jaune]{(59.5,56.2)}{(91,59.6)}
  \surligner[jaune]{(10.5,46.4)}{(50,49.8)}
  \ecrire[vert][9]{(96,61.5)}{$\sqrt{5{,}9174} = 2{,}433 = z_{obs}$}
  \ecrire[vert][9]{(53,49)}{$\leftarrow$ = NOTRE IC À LA MAIN}
}

\annot{56}{QUIZ}{
  \ecrire[violet][9]{(18,41.5)}{$\widehat{\text{ET}} = \sqrt{0{,}682 \times 0{,}318\,(\frac{1}{12} + \frac{1}{10})} = 0{,}199$ \ ($\approx 3$ FOIS PLUS GRANDE QU'AVEC 120 ET 110)}
  \ecrire[rouge][9]{(18,28.5)}{AUSSI : $n_1(1-\widehat{P}) = 12 \times 0{,}318 = 3{,}8 < 5$}
}

% --- Synthèse ---------------------------------------------------------------
\annot{58}{Quel test utiliser}{
  \ecrire[vert][8]{(12,41.3)}{EX. 1 (Q1), EX. 2}
  \ecrire[vert][8]{(5,23.5)}{EX. 3,\\ZINC FOND/\\SURFACE}
  \ecrire[vert][8]{(50,9.8)}{EX. 1 (ZINC, Q2)}
  \ecrire[vert][8]{(110,9.8)}{EX. 2 (MERCURE)}
  \ecrire[vert][8]{(125,41.3)}{EX. 4}
}

\annot{59}{QUIZ}{
  \entourer{(30.5,69.7)}{8.3}{2.4}
  \entourer{(25.3,55.6)}{8}{2.4}
  \entourer{(110.8,41.4)}{6}{2.3}
  \entourer{(102.5,27.4)}{11}{2.4}
  \entourer{(121.5,13.4)}{10}{2.3}
}

\annot{60}{Tableau récapitulatif}{
  \ecrire[vert][7.5]{(128,68.7)}{VAR.TEST()}
  \ecrire[vert][7.5]{(128,61.8)}{T.TEST(..., TRUE)}
  \ecrire[vert][7.5]{(128,55.5)}{T.TEST(..., FALSE)}
  \ecrire[vert][7.5]{(128,49.1)}{T.TEST(PAIRED=T)}
  \ecrire[vert][7.5]{(128,42.6)}{PROP.TEST()}
}

\produire{Chapitre_4}
\end{document}
